\documentclass[12pt,a4paper]{article}

\usepackage{amsmath, amsthm, amssymb, amsfonts}
\usepackage{mathrsfs}
\usepackage{hyperref}
\usepackage{geometry}
\usepackage{enumitem}
\usepackage{cleveref}

\geometry{margin=2.5cm}

\newtheorem{theorem}{Theorem}[section]
\newtheorem{proposition}[theorem]{Proposition}
\newtheorem{lemma}[theorem]{Lemma}
\newtheorem{corollary}[theorem]{Corollary}
\newtheorem{definition}[theorem]{Definition}
\newtheorem{remark}[theorem]{Remark}
\newtheorem{assumption}[theorem]{Assumption}
\newtheorem{openproblem}[theorem]{Open Problem}
\newtheorem{conjecture}[theorem]{Conjecture}

\DeclareMathOperator{\tr}{tr}
\DeclareMathOperator{\Im}{Im}
\DeclareMathOperator{\spec}{spec}
\newcommand{\Xc}{X_c}
\newcommand{\Tc}{\mathcal{T}_c}
\newcommand{\Kc}{\mathcal{K}_c}
\newcommand{\Hil}{\mathcal{H}_{(-1,1)}}
\newcommand{\norm}[1]{\left\|#1\right\|}
\newcommand{\inner}[2]{\langle #1,\, #2 \rangle}
\newcommand{\pswf}[1]{\psi_{#1}^{(c)}}
\newcommand{\lam}[1]{\lambda_{#1}^{(c)}}
\newcommand{\Ai}{\operatorname{Ai}}

\title{%
  \textbf{Galerkin Norm Estimates for the Prolate Phase Operator}\\[0.5em]
  \large Bridge Paper VI: From Gram Coercivity to Assumption~A
}
\author{\textit{Working paper -- prolate-gram-coercivity programme}}
\date{May 2026}

\begin{document}
\maketitle

\begin{abstract}
We study the nonlinear phase operator $\mathcal{T}_c$ arising in the
Riemann--Hilbert formulation of PSWF asymptotics.
The analysis has three layers above an unconditional base.

\smallskip\noindent\textbf{Layer~0 (unconditional):}
(i)~$T_{ii}=0$;
(ii)~$|I_{kl}|\leq C_1c^{-1/2}/|k-l|$, uniformly over bulk and
transition zones after decomposition;
(iii)~the bulk oscillation barrier is localised to
$|k-l|\lesssim c^{1/6}$;
(iv)~$\|T\|\geq\tfrac{\kappa}{2}\log c$ is optimal in the class
of non-oscillatory Schur-type bounds.
These constitute the geometry of the obstruction.

\smallskip\noindent\textbf{Layer~1 (Assumption~B-strong):}
$\norm{D\mathcal{T}_c^{(N)}}\leq C(A)\log c$.
Assumption~B splits into B-weak (mechanically derived,
$P_{kl}\lesssim c^{5/3}|k-l|$) and B-strong
($P_{kl}\leq C_2c^{1/2}$, requires phase-sensitive
Airy--WKB coupling, open).

\smallskip\noindent\textbf{Layer~2 (Assumption~B$'$):}
$\norm{D\mathcal{T}_c^{(N)}}\leq C(A)c^{-1/2}\log c\to 0$.
Assumption~B$'$ requires both phase equidistribution
(Conjecture~\ref{conj:phase}) \emph{and} amplitude regularity
(Conjecture~\ref{conj:amplitude}); neither alone suffices.
The No-go Theorem shows Layer~1 is optimal within the
class of non-oscillatory Schur-type bounds;
contraction requires signed oscillatory structure.
The near/far splitting ($T=T^{\rm near}+T^{\rm far}$,
relative bandwidth $c^{-1/6}\to 0$) reduces the global
contraction problem to a band-matrix of shrinking relative bandwidth;
$T^{\rm near}$ has local Toeplitz symbol structure in the WKB
scaling limit, but $\ell^2$-operator convergence to a Toeplitz
limit is not established.
All gaps are precisely isolated as Open Problems.
\end{abstract}

\tableofcontents
\bigskip\hrule\bigskip

%% ============================================================
\section{Introduction and Context}
%% ============================================================

\subsection{Programme overview}

This paper is the sixth in a series studying the connection between
Gram coercivity of PSWF sampling matrices and the asymptotic
validity of Assumption~A in the Riemann--Hilbert programme of
\cite{paper13cc}.
The central object is the spectral density
\[
  \rho(s;\,g_c)
  =\frac{1}{\pi}\Im\frac{d}{ds}\log\det(I-s\Kc),
\]
where $\Kc$ is the finite-rank sinc operator on $(-1,1)$ with
bandwidth $c>0$.
Assumption~A posits the existence and regularity of the
phase function $g_c$.

\subsection{The nonlinear operator $\mathcal{T}_c$}

The phase $g_c$ satisfies
\begin{equation}\label{eq:fixedpoint}
  g=\Tc(g),\qquad
  (\Tc h)(t)
  =\Hil\!\left[-c\int_t^1\arccos(s)\,\rho(s;\,g_h)\,ds\right](t).
\end{equation}
Proving $\|D\Tc^{(N)}\|<1$ yields Assumption~A via the Banach
fixed-point theorem.

\subsection{Why global $L^2$-contraction fails}

$\norm{D\Tc}_{L^2\to L^2}\gtrsim c^{8/3}$
($\norm{R_s}_{\mathfrak{S}_2}^2\sim c^{5/3}$ times Widom prefactor $c$).
Spectral stability lives on finitely many modes.

\subsection{Main results}

\begin{theorem}[Galerkin norm estimate]\label{thm:main}
Let $N=\lfloor Ac^{1/3}\rfloor$,
$\mathcal{T}_c^{(N)}:=P_N\circ\Tc\circ P_N$.

\medskip\noindent\textbf{(a) Norm bound}
(Assumption~\ref{ass:Bstrong} only):
$\norm{D\mathcal{T}_c^{(N)}}\leq C(A)\log c$.
Not a contraction.

\medskip\noindent\textbf{(b) Contraction}
(Assumptions~\ref{ass:Bstrong} and~\ref{ass:Bprime}):
$\norm{D\mathcal{T}_c^{(N)}}\leq C(A)c^{-1/2}\log c\to 0$.
In particular, $\norm{D\mathcal{T}_c^{(N)}}<1$ for all
$c\geq c_0$, where $c_0=c_0(C,A)$ depends only on the
constants in Assumptions~\ref{ass:Bstrong} and~\ref{ass:Bprime}.
\end{theorem}

\begin{theorem}[No-go theorem]\label{thm:nogo}
$N\times N$ matrix $T$ with $T_{ii}=0$,
$|T_{ij}|\leq\kappa/|i-j|$ for $i\neq j$:
\[
  \norm{T}_{\ell^2}\geq\tfrac{\kappa}{2}H_{N-1}\sim\tfrac{\kappa}{6}\log c.
\]
The bound of Theorem~\ref{thm:main}(a) is optimal in this class.
No argument based solely on non-oscillatory \textup{(}Schur-type\textup{)}
bounds $|T_{ij}|\leq\kappa/|i-j|$ can establish
$\norm{T}<1$; signed oscillatory structure is necessary.
\end{theorem}

\begin{proof}
Classical Schur: $\norm{T}\geq\tfrac{1}{2}\sup_i\sum_{j\neq i}|T_{ij}|
\geq\tfrac{\kappa}{2}H_{N-1}$; $H_{N-1}\geq\tfrac{1}{3}\log c+O(1)$.
\end{proof}

\begin{remark}[Scope of the No-go theorem]\label{rem:nogo_scope}
The No-go theorem applies specifically to the class of
non-oscillatory Schur-type arguments, i.e.\ arguments that
use only $|T_{ij}|\leq f(|i-j|)$ for some non-negative $f$.
It does not preclude contraction via Fourier-multiplier or
Calder\'{o}n--Zygmund methods that exploit signed cancellation
--- those are precisely the content of Layer~2
(Assumption~B$'$, Conjectures~\ref{conj:phase}--\ref{conj:amplitude}).
\end{remark}

\begin{remark}[Structural decomposition]\label{rem:decomposition}
Two independent subproblems:
(1)~\textbf{Pointwise decay} (largely resolved):
$|T_{ij}|\leq C/|i-j|$; mechanism: IBP $+$ Airy $+$ SL.
(2)~\textbf{Summation cancellation} (open):
$\sum_j T_{ij}$ with signs; Calder\'{o}n--Zygmund regime.
Theorem~\ref{thm:nogo}: subproblem~(1) cannot resolve~(2).
\end{remark}

\subsection{Proof architecture}
\label{sec:architecture}

\textbf{Layer~0 (unconditional).}
\begin{enumerate}[nosep,label=(\arabic*)]
  \item $T_{ii}=0$ (\S\ref{sec:matrix});
  \item $|I_{kl}|\leq C_1c^{-1/2}/|k-l|$, uniformly over bulk
        and transition zones after decomposition
        (Lemma~\ref{lem:RL});
  \item barrier localised to $|k-l|\lesssim c^{1/6}$
        (Proposition~\ref{prop:bulk_barrier});
  \item $\|T\|\geq\tfrac{\kappa}{2}\log c$ optimal in the
        non-oscillatory Schur class
        (Theorem~\ref{thm:nogo}).
\end{enumerate}
Items~(1)--(3) constitute the geometry of the obstruction:
they identify precisely where and why classical methods fail,
and establish sharp optimality barriers.
Layer~0 does not provide the resolution mechanism;
that is the content of Layers~1--2.
Item~(4) eliminates the entire class of non-oscillatory
Schur-type approaches.

\medskip\noindent\textbf{Layer~1 (Assumption~B-strong, open).}
$P_{kl}\leq C_2c^{1/2}$ $\Rightarrow$
$|T_{ij}|\leq C_1C_2/|i-j|$ $\Rightarrow$
$\norm{T}\leq C(A)\log c$.
B-weak ($P_{kl}\lesssim c^{5/3}|k-l|$) is mechanically derived
but insufficient; B-strong requires phase-sensitive coupling
(Remarks~\ref{rem:Bweak},~\ref{rem:Bstrong_content}).

\medskip\noindent\textbf{Layer~2 (Assumption~B$'$, open).}
Oscillatory cancellation $\Rightarrow$
$\norm{T}\leq Cc^{-1/2}\log c\to 0$.
Requires both Conjecture~\ref{conj:phase} (phase) and
Conjecture~\ref{conj:amplitude} (amplitude);
see Remark~\ref{rem:Bprime_chain}.

\medskip\noindent
The paper's primary contribution is Layer~0:
the unconditional geometry of the obstruction.
Layers~1--2 identify what remains and why it cannot be
obtained by further non-oscillatory analysis.

\subsection{Comparison with standard PSWF theory}
\label{sec:new}

Standard theory (ORX \cite{OsipovRokhlinXiao2013},
Widom \cite{Widom64}, SP \cite{SP61}) provides pointwise
bounds, eigenvalue asymptotics, WKB approximations.
Three results here are not in the standard theory, and are
independent of Assumptions~B and~B$'$:

\begin{enumerate}[nosep,label=(\arabic*)]
  \item \textbf{Barrier localisation}
        (Prop.~\ref{prop:bulk_barrier}):
        second-IBP obstruction confined to $|k-l|\lesssim c^{1/6}$,
        relative width $c^{-1/6}\to 0$.
  \item \textbf{Schur optimality} (Thm~\ref{thm:nogo}):
        $\log c$ sharp in the non-oscillatory Schur class.
  \item \textbf{Near/far decomposition}
        (OP~\ref{op:nearfar}):
        reduction to band-matrix of shrinking relative bandwidth.
        Item~(3) connects Layer~0 to the theory of semiclassical
        band operators with shrinking relative bandwidth $c^{-1/6}\to0$
        --- a setting where pseudodifferential scaling arguments
        are in principle available, though not yet applied here.
\end{enumerate}

%% ============================================================
\section{Functional Analytic Setup}
%% ============================================================

\subsection{The weighted space $X_c$}

\begin{definition}\label{def:Xc}
$w_c(s):=(\lam{n(s)}-\lam{n(s)+1})^2$;
$\Xc:=\{h\in L^2(-1,1):\int h=0\}$,
$\norm{h}_{\Xc}^2:=\int|h|^2w_c\,ds$.
\end{definition}

Regimes: bulk $w_c\asymp c^{-2}(1-s^2)^{-1}$;
transition $w_c\asymp c^{-2/3}$; outer negligible.

\begin{proposition}[$A_2$ control]\label{prop:A2}
$[w_c]_{A_2}\leq C$ uniformly.
Hence $\norm{\Hil}_{\Xc\to\Xc}=O(1)$
(Coifman--Fefferman).
\end{proposition}

\subsection{Galerkin projection}

\begin{definition}\label{def:Galerkin}
$X_c^{(N)}:=\operatorname{span}\{\pswf{0},\ldots,\pswf{N-1}\}$,
$P_N$ orthogonal projection.
By Gram coercivity (Paper~I):
$\sigma_{\min}(G^{(N)})\geq\alpha_N\sim c^{-1/2}$.
\end{definition}

%% ============================================================
\section{Resolvent Analysis}
\label{sec:resolvent}
%% ============================================================

\begin{lemma}[HS norm]\label{lem:HS}
$\norm{R_s}_{\mathfrak{S}_2}^2\sim c^{5/3}$
($\sim c^{1/3}$ transition modes, each $\sim c^{4/3}$).
\end{lemma}

Combined with Widom: $\norm{D\Tc}_{L^2\to L^2}\gtrsim c^{8/3}$.

%% ============================================================
\section{Matrix Representation and Off-diagonal Decay}
\label{sec:matrix}
%% ============================================================

\subsection{Zones and scales}

\begin{assumption}[Galerkin regime]\label{ass:G}
$k,l\leq N=\lfloor Ac^{1/3}\rfloor$.
\end{assumption}

Zones: bulk $\mathcal{Z}_{\rm bulk}=\{|s|\leq1-c^{-2/3}\}$;
transition $\mathcal{Z}_{\rm tz}$, measure $\sim c^{-2/3}$.

\begin{remark}[Scale ordering]\label{rem:ordering}
$c^{-1}\ll c^{-2/3}\ll c^{-1/2}\ll c^{-1/3}\ll c^{-1/6}\ll 1$.
\end{remark}

\begin{table}[h]
\centering\renewcommand{\arraystretch}{1.3}
\begin{tabular}{lllll}
\hline
\textbf{Object}&\textbf{Scale}&\textbf{Zone}
  &\textbf{Strength}&\textbf{Source}\\
\hline
$\omega_{kl}^{\rm tz}$&$c^{2/3}|k-l|$&tz&surrogate freq.&WKB\\
$\chi_k-\chi_l$&$2c|k-l|$&uniform&SL gap&Widom\\
$\int|W_k^l|$&$c^{-1/6}/|k-l|$&combined&intermediate
  &Cor.~\ref{cor:Wronskian}\\
$\|\pswf{k}\|_{L^\infty({\rm tz})}$&$c^{1/6}$&tz&Airy
  &Lem.~\ref{lem:Airy}\\
$\|\pswf{k}\|_{L^\infty({\rm bulk})}$&$c^{-1/4}(1{-}s^2)^{-1/4}$
  &bulk&WKB&ORX\\
$|F_{kl}|$ (tz)&$c^{-1/3}/|k-l|$&tz&stronger
  &Lem.~\ref{lem:sublD}\\
$|F_{kl}|$ (bulk)&$c^{-1/2}(1{-}s^2)^{1/2}/|k-l|$&bulk&limiting
  &Prop.~\ref{prop:moment}(c)\\
$|I_{kl}^{\rm tz}|$&$c^{-2/3}/|k-l|$&tz&stronger
  &Lem.~\ref{lem:sublD}\\
$|I_{kl}^{\rm bulk}|$&$c^{-1/2}/|k-l|$&bulk&limiting
  &Lem.~\ref{lem:sublD}\\
\hline
$|I_{kl}|$&$c^{-1/2}/|k-l|$&\textbf{global (all zones)}
  &\textbf{integral bound}&\textbf{Lem.~\ref{lem:RL}}\\
$P_{kl}^{\rm weak}$&$c^{5/3}|k-l|$&transition
  &\textbf{B-weak (mech.)}&\textbf{Rem.~\ref{rem:Bweak}}\\
$P_{kl}^{\rm strong}$&$\leq C_2c^{1/2}$&Galerkin
  &\textbf{B-strong (open)}&\textbf{Ass.~\ref{ass:Bstrong}}\\
$|T_{ij}|$&$O(1)/|i-j|$&global&\textbf{Schur input}
  &\textbf{Prop.~\ref{prop:Tij}}\\
TZ Schur&$O(c^{-1/6}\log c)\to0$&TZ&irrelevant
  &Lem.~\ref{lem:zone_schur}\\
Bulk Schur&$\sim C\log c$&bulk&dominant
  &Lem.~\ref{lem:zone_schur}\\
Schur row sum&$[\tfrac{\kappa}{2}\log c,C\log c]$&global
  &\textbf{optimal (Layer~1)}&Thms~\ref{thm:main}(a),\ref{thm:nogo}\\
Schur row sum&$\leq Cc^{-1/2}\log c$&global
  &\textbf{contraction (Layer~2)}&Thm~\ref{thm:main}(b)\\
\hline
\end{tabular}
\caption{Scaling table. B-weak is mechanically derived;
B-strong requires phase-sensitive coupling (open).
Layer~1 is optimal in the non-oscillatory Schur class;
Layer~2 requires signed oscillatory cancellation.
The global bound for $|I_{kl}|$ absorbs all zones.}
\label{tab:scales}
\end{table}

\subsection{Matrix entries and commutator structure}

\begin{equation}\label{eq:Tij_simple}
  T_{ij} = c\cdot
  \frac{|\lam{j}-\lam{i}|}{(1-\lam{i})(1-\lam{j})}
  \cdot I_{ij} =: P_{ij}\cdot I_{ij},\qquad
  I_{ij}:=\int_{-1}^1\arccos(s)\,\pswf{i}(s)\pswf{j}(s)\,ds.
\end{equation}
\begin{equation}\label{eq:diagzero}
  T_{ii}=0\quad\forall\,i.
\end{equation}

\subsection{Airy approximation and Wronskian bound}

\begin{lemma}[Airy]\label{lem:Airy}
$|\pswf{k}|\lesssim c^{1/6}$ in $\mathcal{Z}_{\rm tz}(k)$.
\emph{(Olver \cite{Olver1974}; ORX \cite{OsipovRokhlinXiao2013})}
\end{lemma}

\begin{corollary}[Wronskian bound]\label{cor:Wronskian}
$W_k^l:=(1-s^2)[\pswf{k}'\pswf{l}-\pswf{k}\pswf{l}']$:
$\int|W_k^l|\lesssim c^{-1/6}/|k-l|$.
\end{corollary}

\begin{proof}
TZ: $c^{1/3}\cdot c^{-2/3}=c^{-1/3}$.
Bulk: $c^{1/2}/(c^{2/3}|k-l|)=c^{-1/6}/|k-l|$; bulk dominates.
\end{proof}

\subsection{SL moment and primitive bounds}

\begin{proposition}\label{prop:moment}
For $k\neq l$, $k,l\leq N$:
\begin{enumerate}[label=(\alph*)]
  \item Parity: $k\equiv l\pmod2\Rightarrow\mu_{kl}=0$.
  \item Moment: $|\mu_{kl}|\lesssim c^{-7/6}/|k-l|^2$.
  \item $F_{kl}(s):=\int_s^1\pswf{k}\pswf{l}\,dt
        =W_k^l(s)/(\chi_k-\chi_l)$,\;
        $|F_{kl}|\lesssim c^{-1/2}(1-s^2)^{1/2}/|k-l|$.
\end{enumerate}
\end{proposition}

\subsection{Off-diagonal integral bounds}

\begin{lemma}[Zone decomposition]\label{lem:sublD}
$|I_{kl}^{\rm tz}|\lesssim c^{-2/3}/|k-l|$,\;
$|I_{kl}^{\rm bulk}|\lesssim c^{-1/2}/|k-l|$,\;
$|I_{kl}|\lesssim c^{-1/2}/|k-l|$.
\end{lemma}

\begin{lemma}[IBP integral bound -- global]\label{lem:RL}
For all $k\neq l$, $k,l\leq N$:
\[
  |I_{kl}|\leq C_1c^{-1/2}/|k-l|,
\]
uniformly over bulk and transition zones after decomposition
\textup{(}i.e.\ the zone-dependent bounds of Lemma~\ref{lem:sublD}
are combined into a single global estimate with constants
independent of the zone structure\textup{)}.
\end{lemma}

\begin{lemma}[Zone Schur]\label{lem:zone_schur}
Under Assumption~\ref{ass:Bstrong}:
$\sum_{j\in{\rm TZ}}|T_{ij}|\leq\tfrac{1}{3}C_2c^{-1/6}\log c\to0$;
$\sum_{j\in{\rm bulk}}|T_{ij}|\leq C_1C_2\log c$.
\end{lemma}

\subsection{Bulk oscillation barrier}
\label{sec:bulk_barrier}

\begin{proposition}[Bulk oscillation barrier]\label{prop:bulk_barrier}
WKB in bulk: $\pswf{k}\sim c^{-1/4}(1-s^2)^{-1/4}\cos\theta_k$.

\textbf{(i)} First IBP: $|I_{kl}^{\rm bulk}|\lesssim c^{-1/2}/|k-l|$.

\textbf{(ii)} Difference-phase frequency:
$\partial_s(\theta_k-\theta_l)\asymp|k-l|(1-s^2)^{-1/2}$.

\textbf{(iii)} Second IBP boundary at $s_*=1-c^{-2/3}$ dominates.

\textbf{(iv)} At $s_*$: ratio (boundary term)/(target $c^{-2/3}/|k-l|$)
$= c^{1/6}/|k-l|$.
Therefore:
$|k-l|\lesssim c^{1/6}$: obstruction active;
$|k-l|\gg c^{1/6}$: second IBP succeeds.
The exponent $1/6$ is independently confirmed by the
three-regime scale balance of Remark~\ref{rem:c16_scale}.

\medskip\noindent
\textbf{The bulk obstruction is localised to $|k-l|\lesssim c^{1/6}$.}
Relative bandwidth $c^{1/6}/N=c^{-1/6}\to0$.
This is a Layer~0 (unconditional) result.
\end{proposition}

\begin{remark}[Origin of the $c^{1/6}$ interaction scale]
\label{rem:c16_scale}
The threshold $|k-l|\sim c^{1/6}$ appearing in
Proposition~\ref{prop:bulk_barrier} is not a free parameter
but arises from the intersection of three independent
asymptotic regimes, each imposing its own characteristic scale.

\medskip\noindent
\textbf{(i) Galerkin spectral scaling.}
The natural dimension of the PSWF transition regime forces
the index truncation $N\sim c^{1/3}$, contributing a factor
$c^{1/3}$ to all index-separation estimates.

\medskip\noindent
\textbf{(ii) Airy transition width.}
The turning point $s_*=1-c^{-2/3}$ defines the spatial
resolution scale of the transition zone as $\Delta s\sim c^{-2/3}$,
contributing a factor $c^{-2/3}$.

\medskip\noindent
\textbf{(iii) Second-IBP amplitude scale.}
After the first integration-by-parts
(Proposition~\ref{prop:bulk_barrier}(i)), the residual
boundary term at $s_*$ carries a factor $c^{1/2}$
from the IBP amplitude, as identified in
Proposition~\ref{prop:bulk_barrier}(iv).

\medskip\noindent
\textbf{Scale balance.}
The crossover threshold is the index separation at which
all three regime-specific scales simultaneously become active:
\[
  |k-l| \;\sim\; c^{1/3} \cdot c^{-2/3} \cdot c^{1/2}
  \;=\; c^{1/6}.
\]
The exponent $\tfrac{1}{6} = \tfrac{1}{3} - \tfrac{2}{3}
+ \tfrac{1}{2}$ is heuristically unique within this
three-regime decomposition: any alternative power $c^\alpha$
with $\alpha\neq\tfrac{1}{6}$ would be inconsistent with
at least one of the three scaling laws.
The detailed verification that this balance reproduces
the IBP crossover condition is given in
Proposition~\ref{prop:bulk_barrier}(iv);
this Remark records only the structural origin of the scales.
\end{remark}

%% ============================================================
\section{Norm Estimate and Contraction}
\label{sec:norm}
%% ============================================================

\subsection{Prefactor: B-weak and B-strong}

\begin{remark}[B-weak: mechanically derived]\label{rem:Bweak}
In the transition regime:
$|\lam{k}-\lam{l}|\asymp c^{-2/3}|k-l|$
and $(1-\lam{k})(1-\lam{l})\asymp c^{-4/3}$,
giving
\[
  P_{kl}^{\rm weak}
  :=c\cdot\frac{|\lam{l}-\lam{k}|}{(1-\lam{k})(1-\lam{l})}
  \asymp c^{5/3}|k-l|.
\]
This follows from standard Widom transition asymptotics
without any phase input.
Combined with Lemma~\ref{lem:RL}:
$|T_{ij}|\lesssim c^{5/3}|i-j|\cdot c^{-1/2}/|i-j|=c^{7/6}$,
which does \emph{not} give a $1/|i-j|$ Schur bound.
B-weak alone is insufficient for Theorem~\ref{thm:main}(a).
\end{remark}

\begin{assumption}[B-strong: phase-sensitive coupling]
\label{ass:Bstrong}
For all $k\neq l$, $k,l\leq N$:
\[
  P_{kl}\leq C_2c^{1/2}.
\]
\end{assumption}

\begin{remark}[What B-strong requires]
\label{rem:Bstrong_content}
The reduction from $c^{5/3}|k-l|$ to $c^{1/2}$ is a factor
$c^{7/6}/|k-l|$ beyond Remark~\ref{rem:Bweak}.
This cannot come from spectral theory or amplitude bounds alone.
It requires a \emph{nonlinear coupling} between:
eigenvalue spacing $|\lam{k}-\lam{l}|$,
WKB phase at the turning point $s_*$, and
Airy-function amplitude.
The cancellation must arise from the correlation structure of
$\lam{k}$ and $(1-\lam{k})$ as joint functions of $k$ via
Airy asymptotics.
Concretely, standard methods fail at two points:
\begin{enumerate}[nosep,label=(\alph*)]
  \item \textbf{Non-uniform stationary phase:}
        stationary phase methods do not apply uniformly
        across the transition zone (phase is not globally
        WKB-stable at $s_*$);
  \item \textbf{Correlated Airy factors:}
        the eigenvalue spacing $|\lam{k}-\lam{l}|$
        and the amplitude $(1-\lam{k})(1-\lam{l})$
        are not asymptotically independent ---
        they correlate through the same Airy function,
        and it is precisely this correlation that B-strong
        must exploit.
\end{enumerate}
This is the precise content of Open Problem~\ref{op:prefactor}.

\medskip\noindent
\textbf{B-strong as independent intermediate hypothesis.}
Assumption~B-strong is treated throughout as an independent
intermediate hypothesis: its role in the implication chain
(B-strong $\Rightarrow$ $|T_{ij}|\leq C/|i-j|$ $\Rightarrow$
$\norm{T}\leq C\log c$) is logically self-contained,
and does not presuppose any part of Assumption~B$'$.

\medskip\noindent
\textbf{Epistemic note on shared asymptotic data.}
Any proof of $P_{kl}\leq C_2c^{1/2}$ is expected to involve
the Airy-phase data at $s_*$, which is the same data
that drives the oscillatory cancellation in B$'$.
This does not create a logical dependency
($T_{ij}=P_{ij}\cdot I_{ij}$ is an exact product,
and the implication chain B-strong $+$ B$'$ $\Rightarrow$ Layer~2
is unaffected),
but it suggests that the two assumptions may share
parts of their eventual proof structure.
Although its eventual proof may rely on structural ingredients
overlapping with B$'$, B-strong is formulated as an
independent assumption and its logical role in Layer~1
stands regardless of the proof strategy chosen.
For the programme, this means the conceptually cleanest
path may be to prove B$'$ directly
(via Conjectures~\ref{conj:phase}--\ref{conj:amplitude})
rather than factoring through B-strong as a separate step;
but B-strong remains a valid and useful intermediate target.
\end{remark}

\subsection{Assumption~B$'$ and the two-conjecture structure}

\begin{assumption}[Schur-summability]\label{ass:Bprime}
\[
  \sup_{i\leq N}\sum_{j\neq i}P_{ij}\cdot|I_{ij}|
  \leq C\,c^{-1/2}\log c.
\]
Not implied by Assumption~\ref{ass:Bstrong};
requires signed oscillatory cancellation.
\end{assumption}

\begin{conjecture}[Linear phase equidistribution]
\label{conj:phase}
The WKB phases $\Theta_k^{(c)}$ at $s_*=1-c^{-2/3}$ satisfy,
for $k\leq N$:
\[
  \Theta_k^{(c)}=\alpha^{(c)}k+\beta^{(c)}+O(k^2/c),
  \qquad
  \alpha^{(c)}\in(0,2\pi)\text{ away from }0,\pi.
\]
\end{conjecture}

\begin{conjecture}[Amplitude regularity on dyadic blocks]
\label{conj:amplitude}
For the amplitudes $A_{ij}:=P_{ij}\cdot c^{1/2}$ (normalised),
on each dyadic block $\mathcal{B}_m=\{j:2^m\leq|i-j|<2^{m+1}\}$:
\[
  |A_{ij}-A_{i,j+1}|=o(A_{ij}/2^m)\quad
  \text{uniformly in }i,j\in\mathcal{B}_m,\,m\leq M,
\]
with constants independent of $c$.

The conjecture splits into two regimes with different analytical content:
\begin{enumerate}[nosep,label=(\alph*)]
  \item \textbf{Right regime} ($x_j > 0$, $\lambda_j \approx 0$):
        $A_{ij} \approx c^{1/2}\lambda_i/(1-\lambda_i) = O(1)$,
        nearly constant in $j$.
        Numerically: $\mathrm{TV}(A_{i,\cdot}|_{x_j>0})/A_{\mathrm{avg}}
        \approx 0.63$, consistent with the conjecture;
        Abel summation applies directly.
  \item \textbf{Left regime} ($x_j < 0$, $\lambda_j \sim 1/2$):
        $A_{ij}$ grows monotonically by a factor $\sim 15\times$
        across the Airy window.
        Numerically: $\mathrm{TV}(A_{i,\cdot}|_{x_j<0})/A_{\mathrm{avg}}
        \approx 1.64$ at finite $c$;
        $\mathrm{TV}/(n\cdot A_{\mathrm{avg}}) \sim c^{-0.30}$,
        so the condition holds only asymptotically and slowly.
        Because $A_{ij}$ is \emph{monotone} on this side,
        summation-by-parts with monotone weights (Abel--Dirichlet)
        gives $|S_{\mathrm{left}}| \leq A_{i,\mathrm{far}}
        \cdot |\sum_k \cos(\alpha k)/k|$,
        reducing to the same $O(c^{-1/2}\log c)$ scale provided
        $\alpha$ is bounded away from $0,\pi$
        (Conjecture~\ref{conj:phase}).
\end{enumerate}
\end{conjecture}

\begin{remark}[Inhomogeneity of the amplitude condition]
\label{rem:amplitude_inhom}
On the right side ($x_j>0$), $(1-\lambda_j)\approx 1$ and $A_{ij}$
is flat in $j$: standard Abel summation and the identity
\[
  \sum_{k=1}^N \frac{\cos(\alpha k)}{k}
  = -\log\!\bigl(2|\sin(\alpha/2)|\bigr) + O(N^{-1})
\]
give the full $c^{-1/2}\log c$ gain for this half of each row.
On the left side ($x_j<0$), the uniform TV condition fails at
finite $c$, but monotonicity makes the Abel--Dirichlet theorem
available as a replacement: the signed sum
$S_{\mathrm{left}} = \sum_{k<0} A_{i,i+k}\cos(\alpha k)/|k|$
is controlled by $A_{i,\mathrm{far}} \cdot O(\log N)$, which
after the $c^{-1/2}$ factor from $I_{ij}$ yields $O(c^{-1/2}\log c)$.
The two sides therefore require different summation arguments
but yield the same asymptotic bound.
\end{remark}

\begin{remark}[Two-conjecture structure of B$'$]
\label{rem:Bprime_chain}
Assumption~B$'$ requires both
Conjecture~\ref{conj:phase} and Conjecture~\ref{conj:amplitude}.
Neither alone is sufficient:
\begin{itemize}[nosep]
  \item \textbf{Phase only} (Conj.~\ref{conj:phase}):
        the signed block sum is
        $S_m(i)=\sum_{j\in\mathcal{B}_m}A_{ij}e^{i\alpha(i-j)}/|i-j|$.
        Abel summation gives $|S_m|\lesssim\sup|A_{ij}|/2^m$
        only if $A_{ij}$ is slowly varying in $j$ on scale $2^m$.
        Without amplitude control, no cancellation is guaranteed.
  \item \textbf{Amplitude only} (Conj.~\ref{conj:amplitude}):
        without phase equidistribution, no oscillatory
        cancellation occurs.
\end{itemize}
Together: Dirichlet-kernel cancellation applies blockwise
(the dyadic decomposition is understood in the standard
semiclassical sense, with amplitude regularity uniform in $i$
up to constants independent of $c$, on each block $\mathcal{B}_m$;
this uniformity is implicit in all blockwise Abel summations),
giving $|S_m(i)|\ll 2C_1C_2$ and
$\sum_m|S_m(i)|\leq Cc^{-1/2}\log c$.
The full implication chain is:
\[
  \text{Conj.~\ref{conj:phase}}
  +\text{Conj.~\ref{conj:amplitude}}
  \;\Rightarrow\;
  \text{B$'$}
  \;\Rightarrow\;
  \norm{T}\leq Cc^{-1/2}\log c
  \;\Rightarrow\;
  \text{Assumption~A}.
\]
No arrow reverses in general.
Conjectures~\ref{conj:phase} and~\ref{conj:amplitude} are
\emph{logically independent}: neither implies the other,
and they address distinct aspects of the asymptotic structure
(phase arithmetic at $s_*$ vs.\ amplitude regularity in the Galerkin
index); they are, however, expected to be jointly necessary for B$'$,
since the cancellation mechanism requires both.
The Conjectures are formulated at the level of local WKB data
at the single point $s_*$, and are strictly weaker than
Assumption~B$'$.
They do not resolve the contraction problem ---
they identify its minimal local input.
\end{remark}

\begin{remark}[Schur bound vs.\ spectral norm: rank structure of $T$]
\label{rem:rank}
Assumption~B$'$ is formulated as a Schur row-sum condition
$\sup_i \sum_j P_{ij}|I_{ij}| \leq Cc^{-1/2}\log c$.
Numerically, the signed row sums already satisfy this bound
at finite $c$ (ratio $\sim 1.2$ at $c=100$).
However, the spectral norm satisfies $\|T\|_{\mathrm{op}} \sim c^{0.40}$,
which exceeds $c^{-1/2}\log c$ by a factor growing with $c$.

The explanation is the \emph{low-rank structure} of $T$: numerically,
$T \approx \sigma_1 v_1 v_1^T + \sigma_2 v_2 v_2^T$
with $\sigma_1 \sim c^{0.40}$, $\sigma_2 \ll \sigma_1$,
and all remaining singular values $< 10^{-3}$.
For such a low-rank operator the Schur bound
$\|T\| \leq \max_i \sum_j |T_{ij}|$
is not tight when $v_1$ is far from the constant vector $\mathbf{1}$:
the signed cancellation that makes row sums small
is \emph{not} the same cancellation that suppresses $\|Tv_1\|$.

Therefore Assumption~B$'$ (row sums) is \emph{necessary but not
sufficient} for $\|T\|_{\mathrm{op}} \leq Cc^{-1/2}\log c$
unless the dominant eigenvector satisfies a delocalization condition.
See Open Problem~\ref{op:eigenvector}.
\end{remark}

\subsection{Matrix entry bound and proof of Theorem~\ref{thm:main}}

\begin{proposition}[Matrix entry bound]\label{prop:Tij}
Under Assumption~\ref{ass:Bstrong}: $|T_{ij}|\leq C_1C_2/|i-j|$
for all $i\neq j$, $i,j\leq N$.
\end{proposition}

\begin{proof}
By definition~\eqref{eq:Tij_simple}, $T_{ij}=P_{ij}\cdot I_{ij}$.
Both bounds hold uniformly for $k,l\leq N$ in the Galerkin
regime (Assumption~\ref{ass:G}), and are regime-coherent:
Assumption~\ref{ass:Bstrong} gives $P_{ij}\leq C_2c^{1/2}$
(globally in the Galerkin regime);
Lemma~\ref{lem:RL} gives $|I_{ij}|\leq C_1c^{-1/2}/|i-j|$
(globally, all zones absorbed).
Since $T_{ij}$ is an exact product of these two factors
(not an asymptotic coupling),
$|T_{ij}|\leq C_2c^{1/2}\cdot C_1c^{-1/2}/|i-j|=C_1C_2/|i-j|$.
\end{proof}

\begin{remark}[$c$-power balance]\label{rem:balance}
$c^{+1/2}_{\rm B\text{-}strong}\times c^{-1/2}_{\rm Lem.\ref{lem:RL}}
=c^0$: conditional on B-strong, the norm grows like $\log c$.
Without B-strong: $c^{5/3}|i-j|\times c^{-1/2}/|i-j|=c^{7/6}$
(B-weak alone gives no Schur bound).
\end{remark}

\begin{proof}[Proof of Theorem~\ref{thm:main}]
$T_{ii}=0$ by~\eqref{eq:diagzero}.
\textbf{(a):} Schur $+$ Prop.~\ref{prop:Tij}:
$\norm{T}\leq C_1C_2H_N=:C(A)\log c$.
\textbf{(b):} Assumption~\ref{ass:Bprime}
$\Rightarrow\sup_i\sum_{j\neq i}P_{ij}|I_{ij}|
\leq Cc^{-1/2}\log c\to0$;
the contraction threshold $c_0=c_0(C,A)$ is the smallest
$c$ with $C(A)c^{-1/2}\log c<1$.
\end{proof}

\begin{table}[h]
\centering\renewcommand{\arraystretch}{1.3}
\begin{tabular}{llll}
\hline
\textbf{Step}&\textbf{Bound}&\textbf{Source}&\textbf{Layer}\\
\hline
$T_{ii}=0$&---&commutator&0\\
$|I_{ij}|\leq C_1c^{-1/2}/|i-j|$ (all zones)&---&Lem.~\ref{lem:RL}&0\\
Barrier$\subseteq\{|k-l|\lesssim c^{1/6}\}$&---
  &Prop.~\ref{prop:bulk_barrier}&0\\
$\norm{T}\geq\tfrac{\kappa}{2}\log c$ (non-osc.\ class)&optimal
  &Thm~\ref{thm:nogo}&0\\
\hline
$P_{kl}\leq C_2c^{1/2}$&B-strong&Ass.~\ref{ass:Bstrong}&1 (open)\\
$|T_{ij}|\leq C_1C_2/|i-j|$ (regime-coherent)&---&Prop.~\ref{prop:Tij}&1\\
$\norm{T}\leq C\log c$&optimal&Thm~\ref{thm:main}(a)&1\\
\hline
Conj.~\ref{conj:phase} + Conj.~\ref{conj:amplitude}&---&---&2 (open)\\
$\norm{T}\leq Cc^{-1/2}\log c\to0$, $c\geq c_0$&---&Thm~\ref{thm:main}(b)&2\\
\hline
\end{tabular}
\caption{Proof architecture. Layer~0 unconditional (geometry of obstruction,
not resolution mechanism);
Layer~1 conditional on B-strong (bounded-but-not-contracting);
Layer~2 conditional on both Conjectures (genuine oscillatory cancellation).}
\label{tab:mechanisms}
\end{table}

\begin{corollary}[Conditional fixed point]\label{cor:fixedpoint}
Under B-strong and B$'$: there exists $c_0=c_0(C,A)$ such that
for all $c\geq c_0$, there is a unique $g^{(N)}$ with
$\mathcal{T}_c^{(N)}(g^{(N)})=g^{(N)}$;
Picard rate $\leq Cc^{-1/2}\log c<1$.
\end{corollary}

%% ============================================================
\section{Dyadic Decomposition and Routes to Contraction}
\label{sec:dyadic}
%% ============================================================

\subsection{Dyadic blocks}

$\mathcal{B}_m:=\{j:2^m\leq|i-j|<2^{m+1}\}$,
$m=0,\ldots,M=\lfloor\log_2 N\rfloor$.

\begin{proposition}\label{prop:dyadic_block}
Under B-strong:
$\sum_{j\in\mathcal{B}_m}|T_{ij}|\leq2C_1C_2$ each $m$;
$\sum_{j\neq i}|T_{ij}|\leq2C_1C_2(M+1)\sim C(A)\log c$.
\end{proposition}

\subsection{Structural form of B$'$}

$S_m(i):=\sum_{j\in\mathcal{B}_m}T_{ij}$ (signed).
B$'$ $\Leftrightarrow$ $\sum_m|S_m(i)|\leq Cc^{-1/2}\log c$
$\Leftrightarrow$ $|S_m(i)|\ll 2C_1C_2$ uniformly.
By Theorem~\ref{thm:nogo}: no non-oscillatory Schur improvement
achieves this.

\subsection{Routes to contraction}

\begin{remark}[Four routes]\label{rem:routes}
\begin{enumerate}[nosep,label=(\roman*)]
  \item \textbf{Airy--WKB at $s_*$} (OP~\ref{op:bulk_refined}):
        cancellation of boundary term for $|k-l|\lesssim c^{1/6}$;
        rate $c^{-1/6}\log c$ under B-strong alone.
        Most immediate target; directly controls $T^{\rm near}$.
  \item \textbf{Sign oscillation in $j$}:
        $T_{ij}$ alternates at period $\sim1$
        $\Rightarrow|S_m|=O(1/2^m)$, $\sum|S_m|=O(1)$.
        Requires both Conjectures (Rem.~\ref{rem:Bprime_chain}).
  \item \textbf{Near/far splitting} (OP~\ref{op:nearfar}):
        $T=T^{\rm near}+T^{\rm far}$,
        near-band width $c^{1/6}$, relative bandwidth $c^{-1/6}\to0$.
        Route~(i) controls $T^{\rm near}$;
        $T^{\rm far}$ free of boundary-term obstruction.
  \item \textbf{Weighted norm} (OP~\ref{op:weightnorm}):
        all natural candidates fail (\S\ref{sec:spectral_weights});
        likely equivalent in difficulty to the contraction itself
        (see Remark~\ref{rem:weightnorm_hardness}).
\end{enumerate}
Routes~(i)+(iii) complementary.
Routes~(ii)+(iii) converge under both Conjectures.
\end{remark}

\subsection{Phase structure and Toeplitz model}
\label{sec:fourier}

\begin{remark}[Fourier-multiplier route]\label{rem:fourier_route}
Under Conjecture~\ref{conj:phase}: $T_{ij}\approx f(i-j)/(i-j)$
with $f$ oscillatory and bounded.
Then $\norm{T}_{\ell^2}=\|\widehat{f/|\cdot|}\|_{\ell^\infty}=O(1)$.
This requires Conjecture~\ref{conj:amplitude} for amplitude
uniformity; without it, $f$ depends on $i$ and $j$ separately.
\end{remark}

\begin{remark}[Toeplitz model: local symbol, not limit operator]
\label{rem:toeplitz}
Under Conjecture~\ref{conj:phase},
$T_{ij}\approx g(i-j)$ holds in the near-band
$|i-j|\leq c^{1/6}$.
In this regime $T^{\rm near}$ has \emph{local Toeplitz symbol
structure} in the WKB scaling limit $c\to\infty$,
$k/N\to\xi\in[0,1]$ fixed.
This local symbol structure should not be interpreted as a
global Toeplitz limit; it is a microlocal approximation valid
only within the near-band scaling window $|i-j|\leq c^{1/6}$,
and does not imply operator convergence to a Toeplitz limit
in $\ell^2$.

\medskip\noindent\textbf{Two open points.}
\begin{enumerate}[nosep,label=(\arabic*)]
  \item \textbf{Homogenisation:}
        the error $\|T_{ij}-g(i-j)\|$ in the near-band
        is not controlled; uniform translation-invariance
        requires Conjecture~\ref{conj:amplitude}.
  \item \textbf{Operator convergence:}
        $\ell^2$-convergence of $T^{\rm near}$ to a Toeplitz
        operator is not established
        (OP~\ref{op:nearfar}, part~(2b)).
\end{enumerate}
Conditional on both Conjectures and if $\ell^2$-operator convergence
holds (open point~(2)), the Szeg\H{o}--Toeplitz theorem
\emph{would give}
$\norm{T^{\rm near}}_{\ell^2}\to\norm{g}_{L^\infty(\mathbb{T})}$,
reducing the contraction question to $\norm{g}_{L^\infty}<1$
--- a finite, computable spectral condition.
\end{remark}

\subsection{Why weights fail}
\label{sec:spectral_weights}

\begin{proposition}\label{prop:spectral_weight_fail}
For $k\leq N$: $1-\lambda_k\asymp c^{-2/3}$ uniformly.
Spectral weights $w_k=(1-\lambda_k)^\alpha$:
$w_j/w_i=1+O(c^{-1/3})$, no modification of Schur kernel.
Power weights $(1+k)^\beta$ and exponential $e^{\gamma k/N}$:
row sum $\sim\log N$ for all $\beta,\gamma$.
\end{proposition}

\begin{remark}[Weighted norm: likely equivalent difficulty]
\label{rem:weightnorm_hardness}
A successful two-weight Schur condition requires the weights
to encode the oscillatory cancellation structure of $T_{ij}$
--- information equivalent to controlling the signed sums $S_m(i)$.
Open Problem~\ref{op:weightnorm} is therefore likely no easier
than establishing B$'$ directly.
It is retained as a conceptual route, not a practical shortcut.
\end{remark}

%% ============================================================
\section{Open Problems}
\label{sec:open}
%% ============================================================

\begin{openproblem}[B-strong: phase-sensitive prefactor]
\label{op:prefactor}
Prove Assumption~\ref{ass:Bstrong}: $P_{kl}\leq C_2c^{1/2}$.
The gap is a factor $c^{7/6}/|k-l|$ beyond the mechanically
derived B-weak bound (Remark~\ref{rem:Bweak}),
and arises from the two failure modes of Remark~\ref{rem:Bstrong_content}:
non-uniform stationary phase in the transition zone, and
the correlated Airy-function structure of eigenvalue spacing
and amplitude.
This requires: explicit asymptotic formula for $P_{kl}$ with
controlled error; phase-sensitive Airy--WKB matching at $s_*$;
and control of the joint Airy correlation.
Although B-strong is an independent intermediate hypothesis,
its proof is expected to involve structural ingredients
that may overlap with those needed for
Conjectures~\ref{conj:phase}--\ref{conj:amplitude};
it may therefore be more efficient to prove B$'$ directly.
\end{openproblem}

\begin{openproblem}[Airy--WKB matching at $s_*$]
\label{op:bulk_refined}
Establish cancellation of the boundary term for
$|k-l|\lesssim c^{1/6}$ via Airy connection formulae.
Substitute $s=1-c^{-2/3}x$: difference phase linearises,
boundary term reduces to explicit $\Ai\cdot\Ai$ integral.
If cancellation holds: unconditional rate $c^{-1/6}\log c$
under B-strong alone; simultaneously resolves OP~\ref{op:nearfar}(2a).
\end{openproblem}

\begin{openproblem}[Near/far splitting]\label{op:nearfar}
$T=T^{\rm near}+T^{\rm far}$,
$T^{\rm near}_{ij}:=T_{ij}\cdot\mathbf{1}_{|i-j|\leq c^{1/6}}$.
\begin{enumerate}[nosep,label=(\arabic*)]
  \item Unconditional bound on $\norm{T^{\rm far}}$
        (no boundary-term obstruction for $|i-j|\gg c^{1/6}$).
  \item Control $\norm{T^{\rm near}}$:
        \begin{enumerate}[nosep,label=(\alph*)]
          \item Airy--WKB matching (OP~\ref{op:bulk_refined}).
          \item Homogenisation: prove
                $\|T_{ij}-g(i-j)\|_{\rm op}=o(1)$ in near-band;
                requires Conjecture~\ref{conj:amplitude}
                and $\ell^2$-operator convergence to Toeplitz limit
                (see Remark~\ref{rem:toeplitz}).
        \end{enumerate}
\end{enumerate}
\end{openproblem}

\begin{openproblem}[Amplitude regularity]
\label{op:amplitude}
Prove Conjecture~\ref{conj:amplitude}:
$|A_{ij}-A_{i,j+1}|=o(A_{ij}/2^m)$ on $\mathcal{B}_m$,
with constants independent of $c$.
The right regime ($x_j>0$) is numerically consistent with the
conjecture ($\mathrm{TV}/A_{\mathrm{avg}}\approx 0.63$) and
analytically accessible via standard Airy asymptotics.
The left regime ($x_j<0$) satisfies
$\mathrm{TV}/(n\cdot A_{\mathrm{avg}})\sim c^{-0.30}\to 0$
asymptotically, but the convergence is slow and the uniform
$o(1)$ condition requires explicit control of the Airy-quotient
gradient $\partial_j A_{ij}$ near $\lambda_j\sim 1/2$.
Together with Conjecture~\ref{conj:phase} it implies B$'$.
\end{openproblem}

\begin{openproblem}[Dominant eigenvector delocalization]
\label{op:eigenvector}
Numerically, $T$ is approximately rank~$2$ with leading singular
value $\sigma_1\sim c^{0.40}$ and dominant eigenvector $v_1$
satisfying $\langle v_1,\mathbf{1}\rangle\approx 0$.
Assumption~B$'$ controls $\|T\cdot\mathbf{1}\|$ but not $\sigma_1$,
because the Schur bound is not tight for low-rank operators
when $v_1\not\approx\mathbf{1}$ (Remark~\ref{rem:rank}).
Prove one of:
\begin{enumerate}[nosep,label=(\alph*)]
  \item \textbf{Delocalization:}
        $\|v_1\|_\infty\leq C/\sqrt{N}$,
        so $v_1^T T v_1\leq\|T\|_F/N\to 0$.
  \item \textbf{Left-support bound:}
        $\sum_{i:\,x_i<0}|v_1(i)|^2=o(1)$,
        i.e.\ $v_1$ does not concentrate on the left Airy flank
        where $A_{ij}$ is large.
  \item \textbf{Explicit eigenvector approximation:}
        Identify $v_1$ asymptotically as an Airy-weighted mode
        and verify $v_1^T T v_1\leq Cc^{-1/2}\log c$ directly.
\end{enumerate}
Any of (a)--(c) closes the gap between the row-sum
condition~B$'$ and $\|T\|_{\mathrm{op}}\leq Cc^{-1/2}\log c$.
Conditional on Toeplitz convergence
(Remark~\ref{rem:toeplitz} and OP~\ref{op:nearfar}(2b)),
(a) follows from Fourier delocalization of Toeplitz eigenvectors,
providing a further conditional link between
OP~\ref{op:nearfar}(2b) and OP~\ref{op:eigenvector}.
\end{openproblem}

\begin{openproblem}[Weighted norm]\label{op:weightnorm}
Find $c$-dependent $\norm{\cdot}_w$ with
$\norm{D\mathcal{T}_c^{(N)}}_w<1$ unconditionally.
By Remark~\ref{rem:weightnorm_hardness}, this is likely
equivalent in difficulty to proving B$'$ directly.
A successful weight must encode oscillatory cancellation
and satisfy the two-weight Schur condition.
\end{openproblem}

\begin{openproblem}[Optimal Galerkin dimension]\label{op:Nopt}
Maximal $N(c)$ with Schur bound $<1$.
Conjecture: $N_{\rm opt}\sim c^{1/2}$.
\end{openproblem}

%% ============================================================
\section{Summary and Logical Map}
\label{sec:summary}
%% ============================================================

\subsection*{Layer~0: geometry of the obstruction (primary contribution)}
$T_{ii}=0$;\;
$|I_{kl}|\leq C_1c^{-1/2}/|k-l|$ (all zones, uniform);
barrier $\subseteq\{|k-l|\lesssim c^{1/6}\}$;\;
$\log c$ optimal in non-oscillatory Schur class;
TZ contribution $o(1)$;\;
all natural weights fail.
Layer~0 is sharp, unconditional, and identifies the failure
surfaces of all classical approaches.
It does not provide the resolution mechanism.
The $c^{1/6}$ interaction scale is not a free parameter but
arises from the intersection of three independent asymptotic
regimes (Remark~\ref{rem:c16_scale}); the exponent $1/6$
is heuristically unique within this three-regime decomposition.
The near/far band structure connects Layer~0 to semiclassical
band operator theory (shrinking relative bandwidth $c^{-1/6}\to0$),
a setting where pseudodifferential scaling is in principle available.

\subsection*{Layer~1: conditional on B-strong}
$|T_{ij}|\leq C_1C_2/|i-j|$ (regime-coherent product bound);\;
$\norm{T}\leq C\log c$ (bounded, not contracting).
B-strong is an independent intermediate hypothesis
(Rem.~\ref{rem:Bstrong_content});
its eventual proof may involve structural ingredients
overlapping with B$'$, but its logical role in Layer~1
is self-contained.
B-strong requires phase-sensitive Airy coupling and fails by
two concrete mechanisms (Rem.~\ref{rem:Bstrong_content});
B-weak is insufficient (Rem.~\ref{rem:balance}).

\subsection*{Layer~2: conditional on B-strong $+$ B$'$}
$\norm{T}\leq Cc^{-1/2}\log c\to0$ for $c\geq c_0$;\;
conditional fixed point (Cor.~\ref{cor:fixedpoint}).
B$'$ requires Conj.~\ref{conj:phase} (phase) \emph{and}
Conj.~\ref{conj:amplitude} (amplitude, $c$-uniform);
the two conjectures are logically independent but
jointly necessary for B$'$.
The Schur row-sum condition~B$'$ is necessary but not sufficient
for $\|T\|_{\mathrm{op}}\leq Cc^{-1/2}\log c$ unless the
dominant eigenvector satisfies a delocalization condition
(Rem.~\ref{rem:rank}, OP~\ref{op:eigenvector}).

\subsection*{What is genuinely new}
All Layer~0, all unconditional:
\textbf{(1)} barrier localisation (relative width $c^{-1/6}\to0$);
\textbf{(2)} $\log c$ Schur optimality in non-oscillatory class;
\textbf{(3)} near/far band decomposition;
\textbf{(4)} three-regime scale balance identifying the $c^{1/6}$
crossover as heuristically unique within the given decomposition
(Remark~\ref{rem:c16_scale}).

\subsection*{Core diagnosis}
\textit{Local WKB asymptotics do not determine global operator
cancellation.}
Layer~0 provides the complete geometry of the obstruction:
it is sharp, unconditional, and identifies the failure surfaces
of all classical approaches.
The resolution mechanism --- oscillatory cancellation in the
near-diagonal regime --- is the content of Layers~1--2
and remains open.

%% ============================================================
\section*{Acknowledgements}
%% ============================================================

Part of the \emph{prolate-gram-coercivity} programme.
The B-weak/B-strong split, the two-conjecture structure of B$'$,
the identification of amplitude regularity (Conj.~\ref{conj:amplitude})
as a separate gap, the Toeplitz local-symbol caveat
(scaling model vs.\ limit operator, Szeg\H{o} as conditional
implication under open operator-convergence), the uniformity
clauses in Lemma~\ref{lem:RL} and Proposition~\ref{prop:Tij},
the B-strong failure-mode analysis, the hardness remark on
weighted norms, the explicit contraction threshold $c_0$,
the epistemic coupling remark on B-strong vs.\ B$'$,
the semiclassical band operator anchor,
the three-regime scale balance for the $c^{1/6}$ crossover
(Remark~\ref{rem:c16_scale}, heuristically unique within the
given decomposition),
the inhomogeneous left/right splitting of Conjecture~\ref{conj:amplitude}
with the Abel--Dirichlet argument for the left regime
(Remark~\ref{rem:amplitude_inhom}),
the rank-structure gap between B$'$ and $\|T\|_{\mathrm{op}}$
(Remark~\ref{rem:rank}),
and the dominant eigenvector delocalization problem
(Open Problem~\ref{op:eigenvector})
were developed in working sessions, May 2026.

%% ============================================================
\begin{thebibliography}{99}

\bibitem{paper13cc}
  \textit{Companion paper on RHP asymptotics and Assumption~A},
  prolate-gram-coercivity programme, 2025--2026.

\bibitem{Widom64}
  H.~Widom,
  \textit{Asymptotic behavior of the eigenvalues of certain integral
  equations~II},
  Arch.\ Rational Mech.\ Anal.\ \textbf{17} (1964), 215--229.

\bibitem{HMW73}
  R.~Hunt, B.~Muckenhoupt, R.~Wheeden,
  \textit{Weighted norm inequalities for the conjugate function
  and Hilbert transform},
  Trans.\ Amer.\ Math.\ Soc.\ \textbf{176} (1973), 227--251.

  \bibitem{CF74}
  R.~R.~Coifman, C.~Fefferman,
  \textit{Weighted norm inequalities for maximal functions and
  singular integrals},
  Studia Math.\ \textbf{51} (1974), 241--250.

\bibitem{Kato66}
  T.~Kato,
  \textit{Perturbation Theory for Linear Operators},
  Springer, 1966.

\bibitem{SP61}
  D.~Slepian, H.~O.~Pollak,
  \textit{Prolate spheroidal wave functions, Fourier analysis and
  uncertainty~I},
  Bell System Tech.\ J.\ \textbf{40} (1961), 43--63.

\bibitem{OsipovRokhlinXiao2013}
  A.~Osipov, V.~Rokhlin, H.~Xiao,
  \textit{Prolate Spheroidal Wave Functions of Order Zero},
  Springer, 2013.

\bibitem{Olver1974}
  F.~W.~J.~Olver,
  \textit{Asymptotics and Special Functions},
  Academic Press, 1974.

\bibitem{GarciaTyler2020}
  S.~R.~Garcia, R.~A.~Horn,
  \textit{A Second Course in Linear Algebra},
  Cambridge University Press, 2017.

\end{thebibliography}

\end{document}
